\chapter{\IfLanguageName{dutch}{Stand van zaken}{State of the art}}%
\label{ch:stand-van-zaken}

% Tip: Begin elk hoofdstuk met een paragraaf inleiding die beschrijft hoe
% dit hoofdstuk past binnen het geheel van de bachelorproef. Geef in het
% bijzonder aan wat de link is met het vorige en volgende hoofdstuk.

% Pas na deze inleidende paragraaf komt de eerste sectiehoofding.

\section{Navigatieproblemen op begraafplaatsen}

Wanneer men een bezoek brengt aan een kerkhof, staat rust en reflectie centraal. In de praktijk wordt deze ervaring echter vaak verstoord door bepaalde moeilijkheden, zoals het terugvinden van een specifiek graf. Vooral op grote of complex ingedeelde begraafplaatsen kan dit een tijdrovend en frustrerend proces zijn. Een voorbeeld hiervan is Hillig Meer, waar bezoekers moeite hebben met het lokaliseren van graven, doordat fysieke grafmarkeringen ontbreken. Een bezoeker beschreef deze ervaring als volgt: “Ik stond er gewoon naast. Echt… ik keek er overheen. Het was even verwarrend” \autocite{Meer2025}.

Dergelijke ervaringen tonen aan dat het huidige systeem van fysieke oriëntatie op begraafplaatsen niet altijd volstaat. Het zoeken naar een graf kan leiden tot stress, onzekerheid en een gevoel van ongemak, wat haaks staat op het doel van een bezoek aan een dierbare overledene. Onderzoek van \Autocite{Haekkilae2022} bevestigt dat deze moeilijkheden een negatieve impact hebben op de gebruikerservaring van bezoekers.

\section{Digitalisatie van begraafplaatsen}

In het kader van de bredere digitalisering van de samenleving, worden ook begraafplaatsen steeds vaker uitgerust met digitale hulpmiddelen. Deze evolutie past binnen een algemene trend waarbij fysieke ruimtes ondersteund worden door digitale informatiesystemen om navigatie en toegankelijkheid te verbeteren. Binnen begraafplaatsen vertaalt zich dit voornamelijk naar digitale kaarten, databanken en webapplicaties die informatie over graven centraliseren en toegankelijk maken.

Uit onderzoek van \Autocite{Haekkilae2022} blijkt dat dergelijke digitale toepassingen een duidelijke meerwaarde bieden. Gebruikers geven aan dat het terugvinden van graven aanzienlijk eenvoudiger wordt. Zo vermeldde een deelnemer: “We hadden vijf jaar nodig om het graf van ons familielid te vinden. Het was zeer moeilijk om terug te vinden, zeker in de winter” (P4). Na het gebruik van een digitale applicatie werd deze taak als veel efficiënter en gebruiksvriendelijker ervaren.

\section{Bestaande digitale toepassingen}

Ook in de praktijk bestaan er reeds verschillende digitale oplossingen. Binnen Belgische gemeenten worden bijvoorbeeld systemen gebruikt zoals WebBGP en geoloketten. WebBGP, gebruikt door de gemeente Sint-Gillis-Waas, is een intern systeem dat enkel toegankelijk is voor medewerkers en voornamelijk dient voor beheer en administratie van begraafplaatsen. Hoewel dit systeem efficiënt is voor interne processen, biedt het geen directe ondersteuning aan bezoekers.

Daartegenover staat het geoloket, zoals toegepast in de gemeente Beveren. Dit is een publieke webtoepassing die bezoekers toelaat om via een kaartinterface de locatie van een graf op te zoeken en een route te volgen naar de juiste plaats. Deze oplossing verlaagt de drempel aanzienlijk doordat ze toegankelijk is via een smartphone en geen specifieke technische kennis vereist. Beide systemen werden ontwikkeld door het bedrijf Cevi, wat aantoont dat er reeds gestandaardiseerde oplossingen bestaan binnen dit domein.

\section{Beperkingen van huidige oplossingen}

Ondanks de voordelen van digitalisatie, blijven er belangrijke beperkingen bestaan. Een cruciaal probleem is de gebruiksvriendelijkheid, vooral voor oudere bezoekers. Zoals aangegeven in \Autocite{Haekkilae2022}, ervaren niet alle gebruikers mobiele applicaties als intuïtief. Een gebruiker merkte op: “Zeker oudere mensen gaan het moeilijker hebben, aangezien ze meestal nog geen smartphone gebruiken” (P2).

Daarnaast vereisen huidige oplossingen dat gebruikers hun aandacht verdelen tussen de fysieke omgeving en een scherm. Dit kan leiden tot desoriëntatie of verminderde betrokkenheid bij de omgeving, wat problematisch is in een context waar beleving en respect centraal staan.

\section{Augmented Reality als opkomende technologie}

Om deze beperkingen te overbruggen, wordt in recente literatuur gekeken naar innovatieve technologieën zoals augmented reality (AR). AR integreert digitale informatie rechtstreeks in de fysieke omgeving, waardoor gebruikers geen aparte interface hoeven te raadplegen. Volgens \autocite{Zhao2025} biedt AR een veelbelovende oplossing doordat het de cognitieve belasting vermindert en het situationeel bewustzijn verhoogt. Digitale aanwijzingen kunnen bijvoorbeeld rechtstreeks geprojecteerd worden op het zichtveld van de gebruiker, waardoor navigatie efficiënter wordt.

Studies tonen aan dat mobile augmented reality (MAR) steeds toegankelijker wordt door de brede beschikbaarheid van smartphones en frameworks zoals ARKit en ARCore. \autocite{Xu2018} stelt dat AR-systemen in staat zijn om virtuele objecten in real-time te registreren en uit te lijnen met de fysieke wereld, waardoor een betere gebruikerservaring ontstaat. Deze technologie maakt het mogelijk om tekst, beelden en andere informatie contextueel weer te geven in de omgeving van de gebruiker.

Daarnaast blijkt uit onderzoek dat AR een positieve invloed heeft op de gebruikerservaring. Gebruikers ervaren meer betrokkenheid en een hogere mate van immersie wanneer AR wordt toegepast in navigatie- of ontdekkingscontexten. Dit suggereert dat AR niet enkel functionele voordelen biedt, maar ook de beleving van een bezoek kan verbeteren, wat relevant is in een emotioneel geladen context zoals een begraafplaats.

Tegelijkertijd zijn er ook belangrijke uitdagingen verbonden aan het gebruik van AR in buitenomgevingen. Zo kan het gebruik van AR leiden tot veiligheidsrisico’s, bijvoorbeeld wanneer gebruikers te sterk gefocust zijn op virtuele elementen en minder aandacht hebben voor hun fysieke omgeving. Daarnaast kunnen technische beperkingen, zoals trackingfouten of slechte zichtbaarheid bij fel zonlicht, de gebruikservaring negatief beïnvloeden.

Onderzoek van \autocite{Chen2022} toont bovendien aan dat AR-systemen steeds vaker gecombineerd worden met technieken zoals 3D objectdetectie en -herkenning. Deze systemen maken gebruik van camerabeelden en sensordata om objecten in de omgeving te detecteren en te lokaliseren, waardoor informatie contextueel en dynamisch kan worden aangeboden. Daarnaast kan de integratie met spraakinterfaces bijdragen aan een meer toegankelijke en inclusieve gebruikerservaring.

\section{Link met Proof of Concept}

De literatuur toont aan dat er reeds verschillende digitale oplossingen bestaan voor navigatie op begraafplaatsen, maar dat deze nog beperkingen vertonen op het vlak van gebruiksvriendelijkheid en beleving. Augmented reality wordt hierbij naar voren geschoven als een veelbelovende technologie om deze beperkingen te overbruggen.

Binnen deze bachelorproef wordt deze hypothese concreet onderzocht aan de hand van een proof of concept: een mobiele applicatie die een digitale kaart combineert met een AR-overlay voor navigatie naar een specifiek graf. In tegenstelling tot klassieke geoloketten, die voornamelijk werken met statische kaartweergaven, richt deze applicatie zich op het aanbieden van contextuele en visuele navigatie-instructies rechtstreeks in de omgeving van de gebruiker.

Op basis van de literatuur wordt verwacht dat deze aanpak kan leiden tot een meer intuïtieve, efficiënte en minder stressvolle gebruikerservaring. Tegelijkertijd wordt binnen deze proof of concept ook aandacht besteed aan de eerder geïdentificeerde uitdagingen, zoals gebruiksvriendelijkheid voor diverse doelgroepen, technische betrouwbaarheid en het behoud van een respectvolle gebruikerservaring.




